% ============================================================
% SST Route I theorem target v0.0.4
% Boundary microstates, area density, and relative entropy
% ============================================================

\section{Research-Track boundary-microstate theorem for Route I}
\label{sec:rt_routeI_microstates_v004}

\paragraph{Status.}
\textbf{[CLOSED-CONDITIONAL / MICROSTATE THEOREM].}
The derivation below is exact under the explicitly stated stationary-line,
protected-alphabet, Gaussian-coherent, and reversible-encoding assumptions.
It is not yet a theorem for the complete nonlinear SST substrate.

\subsection{Microstate assumptions}

Let $\Sigma$ be a locally planar boundary element with unit normal
$\widehat{\mathbf n}$. Assume:
\begin{enumerate}
    \item the boundary-relevant line fabric is stationary and ergodic, with
    line-length density $\mathcal L_v$ and tangent distribution
    $f(\widehat{\mathbf t})$;
    \item each independent piercing carries $q\ge2$ protected labels;
    \item weak coherent torsion excitations are represented by product
    equal-covariance Gaussian shifts;
    \item the boundary response is asymptotically reversible and saturates
    the $q$-ary information capacity.
\end{enumerate}

The line-length density is
\begin{align}
    \mathcal L_v
    :=
    \lim_{V\to\infty}
    \frac{\text{total line length in }V}{\operatorname{Vol}(V)},
    \qquad
    [\mathcal L_v]=\mathrm{m^{-2}}.
\end{align}

\subsection{Piercing-density theorem}

A line element $ds$ with tangent $\widehat{\mathbf t}$ has normal projected
measure
\begin{align}
    dA_\perp
    =
    |\widehat{\mathbf t}\cdot\widehat{\mathbf n}|\,ds.
\end{align}
Therefore the crossing density is
\begin{align}
    \boxed{
    n_\perp(\widehat{\mathbf n})
    =
    \mathcal L_v
    \int_{S^2}
    |\widehat{\mathbf t}\cdot\widehat{\mathbf n}|
    f(\widehat{\mathbf t})\,d\Omega
    }.
    \label{eq:rt_v004_nperp_general}
\end{align}
For isotropy, $f=1/(4\pi)$ and
\begin{align}
    \left\langle|\cos\theta|\right\rangle
    =\frac12,
\end{align}
so
\begin{align}
    \boxed{
    n_\perp=\frac{\mathcal L_v}{2}
    }.
    \label{eq:rt_v004_nperp_isotropic}
\end{align}

\subsection{Area-entropy-density theorem}

If each piercing carries $q$ protected labels, then for $N_\perp(A)$
piercings through area $A$,
\begin{align}
    W(A)=q^{N_\perp(A)},
    \qquad
    S_{\rm cap}(A)=N_\perp(A)\ln q.
\end{align}
Ergodicity gives $N_\perp(A)/A\to n_\perp$, hence
\begin{align}
    \boxed{
    \eta_A^{\rm SST}
    =
    n_\perp\ln q
    }.
    \label{eq:rt_v004_eta_general}
\end{align}
Equivalently,
\begin{align}
    \boxed{
    \eta_A^{\rm SST}(\widehat{\mathbf n})
    =
    \mathcal L_v
    \left\langle
    |\widehat{\mathbf t}\cdot\widehat{\mathbf n}|
    \right\rangle_f
    \ln q
    }.
\end{align}
For isotropy,
\begin{align}
    \boxed{
    \eta_A^{\rm SST}
    =
    \frac{\mathcal L_v}{2}\ln q
    }.
    \label{eq:rt_v004_eta_isotropic}
\end{align}
The dimensions are $[\eta_A^{\rm SST}]=\mathrm{m^{-2}}$.

\subsection{Relative entropy from explicit Gaussian boundary microstates}

For each horizon cell $j$, take
\begin{align}
    p_{0,j}(x)
    =
    \frac{e^{-x^2/2}}{\sqrt{2\pi}},
    \qquad
    p_{\phi,j}(x)
    =
    \frac{e^{-(x-\mu_j)^2/2}}{\sqrt{2\pi}}.
\end{align}
Then
\begin{align}
    D_{\rm KL}(p_{\phi,j}\Vert p_{0,j})
    =
    \frac{\mu_j^2}{2}
    \cite{KullbackLeibler1951}.
\end{align}
For midpoint cells $(U_j,\Delta U,\Delta A)$, define
\begin{align}
    \boxed{
    \mu_j
    =
    \sqrt{4\pi(-U_j)\Delta U\Delta A}
    \;\partial_U\widehat\phi(U_j)
    }.
    \label{eq:rt_v004_gaussian_shift}
\end{align}
Additivity of relative entropy gives
\begin{align}
    \delta S_{\rm boundary}^{(N)}
    &:={}
    D_{\rm KL}
    \left(
    \prod_jp_{\phi,j}
    \middle\Vert
    \prod_jp_{0,j}
    \right)
    \\
    &={}
    -2\pi
    \sum_j
    U_j
    \left(\partial_U\widehat\phi(U_j)\right)^2
    \Delta U\Delta A.
\end{align}
Therefore, in the continuum limit,
\begin{align}
    \boxed{
    \delta S_{\rm boundary}
    =
    -2\pi
    \int_{\mathcal H^-}
    U\,T_{UU}\,dU\,dA
    =
    S_{\rm rel}
    }.
    \label{eq:rt_v004_boundary_relative_entropy}
\end{align}

\paragraph{Critical distinction.}
Equal-covariance Gaussian shifts have equal ordinary Shannon entropy, so
\begin{align}
    H[p_{\phi,j}]-H[p_{0,j}]=0
\end{align}
while their relative entropy is positive. Thus
Eq.~\eqref{eq:rt_v004_boundary_relative_entropy} concerns relative
state distinguishability, not a plain absolute-entropy difference. This is
consistent with the relative-entropy formulation of Dorau and Much
\cite{DorauMuch2026}.

\subsection{Reversible encoding and area response}

Quantum/classical Stein asymptotics give relative entropy an operational role
as a distinguishability exponent \cite{HiaiPetz1991}. Under reversible
capacity-saturating $q$-ary encoding, $S_{\rm rel}$ nats activate
\begin{align}
    \delta N
    =
    \frac{S_{\rm rel}}{\ln q}
\end{align}
channels. Hence
\begin{align}
    \delta A
    =
    \frac{\delta N}{n_\perp}
    =
    \frac{S_{\rm rel}}{n_\perp\ln q}.
\end{align}
Using Eq.~\eqref{eq:rt_v004_eta_general},
\begin{align}
    \boxed{
    \eta_A^{\rm SST}\delta A
    =
    \delta S_{\rm boundary}
    =
    S_{\rm rel}
    }.
    \label{eq:rt_v004_encoding_closure}
\end{align}

\subsection{Core-scale no-go theorem}

Let non-overlapping tubes of radius $r_c$ occupy volume fraction
$0<\varphi_v\le1$. Then
\begin{align}
    \varphi_v
    =
    \pi r_c^2\mathcal L_v,
\end{align}
so for isotropy
\begin{align}
    \boxed{
    \eta_A^{\rm SST}
    =
    \frac{\varphi_v\ln q}{2\pi r_c^2}
    }.
    \label{eq:rt_v004_core_eta}
\end{align}
For the strongest minimal binary model,
$\varphi_v=1$, $q=2$, and
$r_c=1.40897017\times10^{-15}\,\mathrm m$,
\begin{align}
    \eta_{A,\max}^{(r_c,q=2)}
    =
    5.557020457583\times10^{28}\,\mathrm{m^{-2}}.
\end{align}
Only as an external post-derivation audit, the observed Newton constant implies
\begin{align}
    \eta_A^{\rm GR}
    =
    \frac{c^3}{4\hbar G}
    =
    9.570182792895\times10^{68}\,\mathrm{m^{-2}}.
\end{align}
Therefore
\begin{align}
    \boxed{
    \frac{\eta_A^{\rm GR}}
    {\eta_{A,\max}^{(r_c,q=2)}}
    =
    1.722178794544\times10^{40}
    }.
    \label{eq:rt_v004_no_go_ratio}
\end{align}
Thus independent binary $r_c$-scale piercings are rejected as a complete
microphysical explanation of the gravitational area coefficient.

\paragraph{Research status.}
The theorem closes the algebraic microstate chain only under its assumptions.
Physical promotion requires independent derivations of $\mathcal L_v$, $q$,
the piercing correlations, and the reversible channel-activation dynamics.
The factor in Eq.~\eqref{eq:rt_v004_no_go_ratio} must emerge without importing
$G$ or $L_p$, or this line-piercing realization of Route I must be abandoned.

% Suggested bibliography entries:
% \bibitem{KullbackLeibler1951} ...
% \bibitem{HiaiPetz1991} ...
% \bibitem{DorauMuch2026} ...
